\begin{abstract}
真实视频里，文本、声音和画面并不总能同时观测，情感判断还需要能回到原始素材上核对。本文把三问当成一条链：构造可反查的特征，在局部缺失下预测情感，再给出可回溯的解释。

针对问题一，我们用 WhisperX 对给定转写做强制对齐得到词级时间，把 BERT 子词、openSMILE 声学帧和 OpenFace 人脸动作单元帧按词时间窗对齐池化，得到768、25、35维特征，每个位置都能反查到原词、秒数和帧号。100条全部保留：77条正常、19条需复核、4条找不到达标的人脸帧；文本100条可用，声学、视觉各92条可用。人工核验的583个词全部通过。

针对问题二，我们把补齐填充、原生不可观测、人为遮蔽三种零值分开记录，建立掩码感知的池化双头模型 PM-B2，同时输出类别和强度。在7种遮蔽组合、3种比例、3种位置构成的63档缺失面板上，三种子平均 macro-F1 为0.5752$\pm$0.0013、MAE 为0.6213$\pm$0.0123，均优于7个基线；测试集完整输入 macro-F1 为0.6147、MAE 为0.6485。遮蔽含文本时 F1 比只遮蔽声学或视觉低约0.04，缺失比例从10\%增至50\%时 F1 由0.5905降到0.5595，位置的影响在验证集上不到0.005。主要错误在中性类，测试集中性召回仅0.361。另做的15种改进配置中，完整输入 F1 最高达0.6244，但在更严格的整模态压力面板上只有0.4139，故未替换正式模型。附件3的30条预测为正向14、中性10、负向6条。

针对问题三，我们固定完整输入下的预测类别，对三模态的8个子集精确计算 Shapley 值，全部样本的效率恒等式误差不超过$10^{-6}$；验证集上文本对分类的平均贡献约为声学、视觉的6～7倍。在附件4上删除得分最高的30\%位置，目标类概率平均下降约0.14，删除随机位置仅约0.03。附件4中以非文本为主导的5条样本置信度都不超过0.565，其余15条平均为0.79。文字与声音不符的样本15，其词时间被人工核验全部拒绝，只给出位置级证据。

\keywords{多模态情感分析；强制对齐；局部模态缺失；池化双头模型；Shapley 归因；删除检验}
\end{abstract}
